\clearpage
\phantomsection% 
\addcontentsline{toc}{chapter}{KATA PENGANTAR}
%\thispagestyle{fancy}

\begin{justifying}
	\large \bfseries \centering \MakeUppercase{Kata Pengantar}\linebreak
	
	\normalsize \normalfont \justifying
	Puji dan syukur penulis panjatkan kepada Tuhan Yesus Kristus atas kasih, penyertaan, dan berkat-Nya yang tidak pernah berhenti dalam kehidupan penulis, sehingga penulis dapat melalui setiap proses dan akhirnya menyelesaikan tugas akhir ini dengan baik.\par
	
	Perjalanan penyusunan tugas akhir ini bukanlah hal yang mudah. Ada banyak proses, tantangan, dan momen jatuh bangun yang harus dilalui. Oleh karena itu, penulis ingin mengucapkan terima kasih kepada:\par
	\begin{enumerate}
		\item Diri penulis sendiri, yang telah bertahan, berjuang, dan tidak menyerah dalam menyelesaikan tugas akhir ini.
		\item Kedua orang tua tercinta, yang selalu memberikan doa, kasih sayang, serta dukungan tanpa henti dalam setiap langkah penulis.
		\item Saudara penulis, Yunus Andreas Sianturi, David Marcellino Sianturi, dan Rifayel Astrid Sianturi, yang selalu memberikan semangat, dukungan, dan menjadi penguat bagi penulis selama proses ini.
		\item Bapak Martin Clinton Tosima Manullang, selaku dosen pembimbing yang dengan sabar membimbing, memberikan arahan, serta masukan yang sangat berarti bagi penulis.
		\item Ibu Leslie Anggraini dan Bapak I Wayan Wiprayoga Wisesa selaku dosen penguji yang telah memberikan kritik dan saran yang membangun demi perbaikan tugas akhir ini.
		\item Teman-teman penulis yang telah hadir dan memberikan semangat dalam proses ini, khususnya teman-teman kontrakan 'BUMN SQUAD', grup 'Game/DL gacor', teman-teman ruang TA Informatika, serta teman-teman lainnya yang tidak dapat disebutkan satu per satu.
	\end{enumerate} \par
	Akhir kata, penulis berharap semoga tugas akhir ini dapat memberikan manfaat dan menjadi langkah awal yang baik untuk perjalanan selanjutnya.
	\vfill
	
\end{justifying}
\clearpage





\begin{comment}
\clearpage
\pagestyle{fancy}
\fancyhf{}
\fancyhead[R]{\thepage}
\phantomsection% 
%\clearpage
%\phantomsection% 
\addcontentsline{toc}{chapter}{KATA PENGANTAR}
%\thispagestyle{fancy}

\begin{justifying}
	\large \bfseries \centering \MakeUppercase{Kata Pengantar}\linebreak
	
	\normalsize \normalfont \justifying
	Puji syukur kehadirat Tuhan Yang Maha Esa atas limpahan rahmat, karunia, serta petunjuk-Nya sehingga penyusunan tugas akhir ini telah terselesaikan dengan baik. Dalam penyusunan tugas akhir ini penulis telah banyak mendapatkan arahan, bantuan, serta dukungan dari berbagai pihak. Oleh karena itu pada kesempatan ini penulis mengucapan terima kasih kepada: \par
	\begin{enumerate}
		\item Bapak Prof. Dr. I. Nyoman Pugeg Aryantha selaku Rektor Institut Teknologi Sumatera.  
		\item Bapak Hadi Teguh Yudistira, S.T., Ph.D. selaku Dekan Fakultas Teknologi Industri.
		\item Bapak Andika Setiawan, S. Kom., M. Cs. selaku Ketua Program Studi Teknik Informatika.
		\item Bapak Ilham Firman Ashari, S. Kom., M.T. selaku Koordinator Tugas Akhir Program Studi Teknik Informatika.  
		\item Bapak Martin C. T. Manullang, Ph.D. selaku Dosen Pembimbing atas ide, waktu, tenaga, perhatian, dan masukan yang telah disumbangsihkan kepada penulis.
            \item Bapak Andika Setiawan, S. Kom., M. Cs dan Bapak Eko Dwi Nugroho, S.Kom., M.Cs. selaku Dosen Penguji atas saran dan masukan yang diberikan. 
		\item Kedua Orang Tua dan Adik yang selalu memberikan dukungan dan doa sehingga penulis dapat menyelesaikan tugas akhir ini. Kelembutan, dukungan, dan cinta yang kalian berikan selalu menjadi sumber inspirasi dan kekuatan.
		\item Teman-teman penulis yang membantu selama masa perkuliahan yang tidak bisa disebutkan satu persatu.
	\end{enumerate} \par
	Akhir kata penulis berharap semoga tugas akhir ini dapat memberikan manfaat bagi kita semua. Penulis menyadari bahwa tugas akhir ini tidak luput dari kekurangan dan kelemahan, dan penulis  terbuka untuk menerima saran, kritik, dan masukan yang membangun.
	\vfill
	
\end{justifying}
\clearpage
\end{comment}
